\begin{afterword}
\summaryhead

The post argues that artificial intelligence is ``fundamentally about reducing the mental to the non-mental,'' and that living inside a mind does not teach this. The brain's work is hidden, so we predict what a ``contemplative'' thing will do only by imagining being contemplative. Aristotle's view that the brain cools the blood, and so makes humans calm, explains nothing; neither does saying that many neurons produce contemplation by ``emergence.'' Both draw arrows between labelled boxes and leave the prediction to empathy.

Real reductions let you say what a system does without pointing at a human: Bayes's theorem for evidence, after Jaynes; Pearl's conditional independence for causal reasoning; the minimax search tree for chess. Designs built on common-sense knowledge, massive parallelism or likeness to the brain are dreams, and so are diagrams of labelled boxes. The author's rule: fund no one who says an AI will work ``just like humans'' until they explain what it does and why. Teaching fundamental AI work therefore means teaching reductionism, Taboo and the avoidance of anthropomorphism.

\responsehead

The post states a sound standard and illustrates it well. An explanation of a mental ability should let you predict it from parts that are not themselves mental. The standard is not new, and the post does not claim it is: Dennett had described AI's progress in 1978 as replacing an intelligent homunculus with a committee of ``relatively ignorant, narrow-minded, blind'' ones. The three examples of real reduction are accurate, the first two are credited to Jaynes and Pearl, and the burglar-alarm case is Pearl's own.

One step goes beyond the evidence: the claim that ``you can only reason about `contemplative-ness' by empathic inference.'' That is simulation theory, one side of an open debate, stated as fact, as it was in ``Humans in Funny Suits.''

The opponents are not named. ``AGI wannabes'' and ``many people'' are unidentified, and no one who holds the listed beliefs is quoted. Two of those beliefs match large programmes of the time, Cyc's hand-entered common sense and connectionism.

Robin Hanson objected at the time that the rule for investors would also refuse whole brain emulation, which copies a brain rather than arguing from analogy to it. The author accepted that: emulation ``might be physically possible'' but was not worth funding. That is a disagreement about emulation, not an oversight.

In short: a sound standard for what counts as explaining intelligence, with accurate examples of real reductions, set against opponents who go unnamed, and a rule for funding that its author, when challenged, applied to brain emulation too.
\end{afterword}
